\section{A quadratic example, worked by hand}\label{ch12:sec:example}
For $T(x)=(x+1)/2$ we knew the fixed point from the start. Now we do not, and the theorem tells us about it before we have found it. Let
\[
T(x)=\frac{x^2+1}{3}\qquad\text{on }X=[0,1],
\]
with the usual distance $|x-y|$. First we check each hypothesis.
\begin{itemize}
\item \emph{The space.} The interval $[0,1]$ is nonempty, and complete by Section~\ref{ch11:sec:complete}.
\item \emph{The self-map condition.} If $0\le x\le1$, then $0\le x^2\le1$, so $1\le x^2+1\le2$ and $1/3\le T(x)\le2/3$. In particular $T(x)\in[0,1]$.
\item \emph{The contraction inequality.} For $x,y\in[0,1]$, factor the difference of two squares and use $|ab|=|a|\,|b|$ and $0\le x+y\le2$:
\[
|T(x)-T(y)|=\frac{|x^2-y^2|}{3}=\frac{|x-y|\,|x+y|}{3}\le\frac23|x-y|.
\]
So $q=2/3$ is a contraction factor.
\end{itemize}
By Theorem~\ref{thm:banach}, $T$ has exactly one fixed point $x_*$ in $[0,1]$, and the iterates from every starting point in $[0,1]$ converge to it.

\subsection*{Iterates and a certificate}
Start from $x_0=0$, and compute with exact fractions, not rounded decimals:
\begin{align*}
x_1&=\frac{0^2+1}{3}=\frac13,\\
x_2&=\frac{(1/3)^2+1}{3}=\frac{1/9+9/9}{3}=\frac{10/9}{3}=\frac{10}{27},\\
x_3&=\frac{(10/27)^2+1}{3}=\frac{100/729+729/729}{3}=\frac{829/729}{3}=\frac{829}{2187}.
\end{align*}
Take $z=x_2$. Then $Tz=x_3$, so the residual of $z$ is $|x_2-x_3|$. Writing $x_2=810/2187$ over the denominator $2187=27\cdot81$,
\[
d(z,Tz)=\left|\frac{810}{2187}-\frac{829}{2187}\right|=\frac{19}{2187}.
\]
Since $1/(1-q)=1/(1-2/3)=3$, the residual estimate \eqref{eq:residual} gives the certificate
\[
d(x_2,x_*)\le3\cdot\frac{19}{2187}=\frac{19}{729}=0.02606\ldots.
\]
A certificate rounded to a decimal must be rounded up, never down; otherwise it is no longer a guarantee. Rounded up, the certificate says that $x_2=10/27\approx0.370$ lies within $0.027$ of the fixed point, although we do not yet know where the fixed point is.

The other two estimates of Section~\ref{ch12:sec:estimates} give weaker bounds. The a priori estimate, with $D=d(x_1,x_0)=1/3$ and $n=2$, uses only the first step:
\[
d(x_2,x_*)\le\frac{(2/3)^2}{1-2/3}\cdot\frac13=\frac49=0.444\ldots.
\]
The a posteriori estimate uses the step from $x_1=9/27$ to $x_2=10/27$ and gives $d(x_2,x_*)\le\frac{2/3}{1/3}\cdot\frac1{27}=\frac2{27}=0.0740\ldots$. The residual estimate is the sharpest, because it uses the extra information in $x_3$.

\subsection*{Checking against the exact solution}
Here we can also find $x_*$ exactly. The fixed-point equation $(x^2+1)/3=x$ becomes $x^2+1=3x$, that is, $x^2-3x+1=0$. Multiply by $4$ and add $5$ to both sides to complete the square:
\[
4x^2-12x+9=5,\qquad\text{that is,}\qquad(2x-3)^2=5.
\]
If $a^2=5$, then $(a-\sqrt5)(a+\sqrt5)=a^2-5=0$, so $a=\pm\sqrt5$. Hence $2x-3=\pm\sqrt5$, and $x=(3\pm\sqrt5)/2$. Since $4<5<9$, we have $2<\sqrt5<3$, so $(3+\sqrt5)/2>5/2$ lies outside $[0,1]$, while $(3-\sqrt5)/2$ lies between $0$ and $1/2$. The fixed point in $[0,1]$ is therefore
\[
x_*=\frac{3-\sqrt5}{2}\approx0.382.
\]
The true error of $x_2$ is $x_*-10/27\approx0.0116$, comfortably below the certificate.

On all of $\RR$, the same formula has the two fixed points $(3\pm\sqrt5)/2$. This does not contradict uniqueness: Step~5 shows that a contraction never has two different fixed points, so $T$ is not a contraction of $\RR$. The uniqueness in Theorem~\ref{thm:banach} is about fixed points that lie in $X$.
